\section{The formal Frobenius sequence: source targets and the cocycle step}
\label{sec:faithful-frobenius-cocycle}

All source references in this section are to the frozen tag
\srcid{alpoge-source-7cbf7a0c}, whose source revision is
\srcid{7cbf7a0c32fca1badf05931879a37c3f8a2b03b4}.
The proof mechanism below is the one in the source: finite jets,
the first partial biextension law, a normal-scheme cocycle root,
and changes of integral formal coordinates.  No identification
with a crystalline or Dieudonne module is used here.

\subsection{The exact targets and their inputs}
\label{subsec:faithful-frobenius-targets}

Write \(O=\mathcal O_E\), with its specified involution.
A formal biextension datum \(B\) of dimension \(d\) consists
of the source's formal group law, its \(O\)-action, and the
two partial multiplication, symmetry and adjoint series,
with the identities called \srcid{FB} in the source.
We retain that datum; the identities used in each proof
are written at their use.
For a one-variable vector series \(f\) and a two-variable
scalar series \(c\), set
\[
 \begin{split}
 \operatorname{lin}_m(f)_{ij}&=[X_j^m]f_i(X),\\
 \operatorname{bil}_{m,n}(c)_{ij}&=[X_i^mY_j^n]c(X,Y).
 \end{split}
\]
Let \(R_a=\operatorname{lin}_1(B.\mathrm{act}(a))\),
\(S=\operatorname{bil}_{1,1}(B.\mathrm{sy})\), and
\(C_a=\operatorname{bil}_{1,1}(B.\mathrm{ad}(a))\).
The source's full formal module is the vector space
\(H_B=k^d\oplus(k^d)^\vee\), in the displayed standard
coordinates, with action
\[
       \rho_H(a)(v,f)=(R_av,R_{\bar a}^{\,t}f-C_a^{\,t}v).
 \tag{FC1}\label{eq:faithful-formal-H-action}
\]
For a prime \(p=\operatorname{char}k\), put
\[
 \begin{split}
 P_B(X,Y)&=\prod_{\nu=0}^{p-1}
       B.\mathrm{ph1}([\nu]X,X,Y),\\
 H&=\operatorname{lin}_p([p]),\qquad
 A=\operatorname{bil}_{p,1}(P_B),\\
 \alpha_B(u)&=(Hu,A^tu),\\
 \beta_B(v,f)&=H^tf-(A-H^tS)v .
 \end{split}
 \tag{FC2}\label{eq:faithful-alpha-beta}
\]
Thus the predicate \srcid{ConjInj B p} means precisely
that the \(k\)-linear map \(\alpha_B\) is injective.
The sign in \(\beta_B\) is the source's sign.

Here are the precise geometric hypotheses in
\srcid{FamilyEx.fib\_inj}.
Let \(E\) be a CM number field with host \(H_E\), and identify
the involution on \(O\) with \(H_E.\mathrm{bar}\).
Let \(t\) be the source's polarized semiabelian \(O\)-family
\(\mathrm{SAB}(H_E,X)\).
The condition \(\mathrm{PowFns}(t)\) says that for every
\(n\ge0\) the pullback
\[
       \Gamma(X,\mathcal O_X)\longrightarrow
                  \Gamma(G^n_X,\mathcal O_{G^n_X})
\]
is bijective.
For a field-valued point \(x:\operatorname{Spec}k\to X\),
\(\mathrm{GeomPF}(t,x)\) is the existence of a torus extension
\[
          1\longrightarrow T\longrightarrow G_x
                     \stackrel q\longrightarrow B_0\longrightarrow1
\]
and a line \(M\) on its abelian quotient, with the source's
\(\mathrm{KTriv}(B_0,M)\), and an isomorphism
\(t.\mathrm{PP}.L|_{G_x}\simeq q^*M\).
The predicate \(\mathrm{FibPF}(t,z)\) asks for such a
presentation at some algebraically closed field-valued
point whose image is \(z\).
These are hypotheses about the actual family; this
section does not manufacture them from semistability.

For every field \(k\) of characteristic \(p\), every
\(x:\operatorname{Spec}k\to X\), and every actual formal
germ \(g\) of \(t\) at \(x\), the source target is
\[
 \mathrm{PowFns}(t)\ \land\
 \mathrm{FibPF}\bigl(t,x(\operatorname{closedPoint}k)\bigr)
                  \quad\Longrightarrow\quad
                         \alpha_{g.B}\text{ is injective}.
 \tag{FC3}\label{eq:faithful-fibre-inj-target}
\]
The prime \(2\) is included.
The germ here contains its actual finite-jet maps to \(G\),
the tangent basis, compatible line frames at all jet orders,
their truncation and axes identities, and the formal
biextension datum \(g.B\).  It is not merely an abstract
formal group equipped with matrices.

The source also uses the intermediate property
\(\mathrm{TanPF}(t,x,p)\).
Put \(R_m=k[\delta]/(\delta^{m+1})\).
For every \(y\in G_x(R_p)\) reducing to the identity,
it requires
\[
 [p]y=0,\qquad
 t_{y_1}^*L\otimes L^{-1}\text{ has a global generator}
                   \quad\Longrightarrow\quad y_1=0,
 \tag{FC4}\label{eq:faithful-tanpf}
\]
where \(y_1\) is its restriction to \(R_1\).
This is exactly the source's translation-difference line
on \(G_{x,R_1}\), not just a line on its formal completion.

For the family conclusion, let \(p\) be prime and let
\(K\) be a field of characteristic \(p\).
The source uses an abstract arithmetic curve \(C\) over
this \(K\): an index set \(I_C\),
integer-valued multiplicative valuations \(v_i\), weights
\(w_i\in\mathbb N\), finite support of \(v_i(a)\ne1\) for
every \(a\ne0\), and the weighted product formula.
Write \(R_i=\{a:v_i(a)\le1\}\).
A lattice bundle is a collection of \(R_i\)-lattices with
local bases, equal to one fixed basis lattice outside a
finite set.  A short exact sequence of such bundles means
both exactness over \(K\) and exactness of the indicated
lattices at every \(i\).
For a curve in the actual family, \srcid{FamilyEx.family\_ex}
constructs generic formal data \(B\), integral data \(B_i\)
over \(R_i\), their residue data, and compatible generic
isomorphisms
\[
                   (\varphi_i,\psi_i,\lambda_i):
                                  B_{i,K}\longrightarrow B.
\]
Here \(\varphi_i,\psi_i\) are inverse formal coordinate
changes and \(\lambda_i\) is the axes-normalized
biextension gauge.  The residue \(\alpha\) is injective,
and the coefficients of their linear and mixed terms
are integral outside a finite set.
Assuming \(I_C\ne\varnothing\), the ensuing target is
\[
 \begin{gathered}
 \mathcal L_i=D\varphi_i(0)\,R_i^d,\qquad
 \mathcal H_i=T_i(R_i^d\oplus R_i^d),\\
 T_i(v,f)=\bigl(P_iv,Q_i^t(f+Z_i^tv)\bigr),\\
 P_i=D\varphi_i(0),\quad Q_i=D\psi_i(0),\quad
 Z_i=\operatorname{bil}_{1,1}(\lambda_i),
 \end{gathered}
 \tag{FC5}\label{eq:faithful-family-lattices}
\]
and the two exact lattice sequences
\[
 \begin{split}
 0&\longrightarrow\mathcal L^\vee
       \stackrel{f\mapsto(0,f)}\longrightarrow\mathcal H
       \stackrel{(v,f)\mapsto v}\longrightarrow\mathcal L
       \longrightarrow0,\\
 0&\longrightarrow F^*\mathcal L
       \stackrel\alpha\longrightarrow\mathcal H
       \stackrel\beta\longrightarrow F^*(\mathcal L^\vee)
       \longrightarrow0.
 \end{split}
 \tag{FC6}\label{eq:faithful-two-sequences-target}
\]
The \(O\)-action on the dual is conjugated, as expressed by
\[
 \rho_H(a)(0,f)=(0,R_{\bar a}^{\,t}f),\quad
 \alpha(R_a^{(p)}u)=\rho_H(a)\alpha(u),\quad
 \beta(\rho_H(a)y)=(R_{\bar a}^{\,t})^{(p)}\beta(y).
 \tag{FC7}\label{eq:faithful-frobenius-equivariance}
\]
Here \(F^*\) is the lattice span of coordinatewise \(p\)-th
powers, not an assumption that \(K\) is perfect.

The final degree target is
\srcid{SplitHermitianGeneratorsExist.deg\_quot\_eq\_zero}.
In addition to (FC6)--(FC7), its hypotheses are:
every integral multiplication characteristic polynomial
of \(O\) splits over \(K\); a character \(\chi:O\to K\);
and, for one fixed \(r\), at every \(i\) a free
\(O\otimes R_i\)-basis of \(\mathcal H_i\).
Precisely, if \(b_1,\ldots,b_{[E:\mathbb Q]}\) is the
specified integral basis of \(O\), there are \(x_1,\ldots,x_r\)
such that all \(\rho_H(b_j)x_l\) are \(K\)-linearly independent
and their \(R_i\)-span is \(\mathcal H_i\).
With the source's quotient lattice, the target is
\[
 \deg_C\!\left(
 \mathcal H\Big/
   \sum_{a\in O}\operatorname{im}(\rho_H(a)-\chi(a))
                    \right)=0.
 \tag{FC8}\label{eq:faithful-character-degree-target}
\]
The local freeness in this statement is a separate input;
neither (FC3) nor (FC6) asserts it.

\subsection{Roots of the additive overlap cocycle}
\label{subsec:faithful-cech-p-root}

\begin{srlemma}\label{lem:faithful-cech-p-root}
Let \(Z\) be an integral normal scheme with function field
\(K_Z\) of characteristic \(p>0\), where \(p\) is prime.
Let \(U_j\) be arbitrary open subsets, and choose
\[
 h_j\in\Gamma(U_j,\mathcal O_Z),\qquad
 n_{jj'}\in\Gamma(U_j\cap U_{j'},\mathcal O_Z)
\]
such that
\[
                         n_{jj'}^p=h_j-h_{j'}.
 \tag{FC9}\label{eq:faithful-p-root-overlap}
\]
Restrictions to intersections are understood.
Suppose some \(U_{j_0}\ne\varnothing\) and some
\(r\in K_Z\) satisfies \(r^p=h_{j_0}\) in the function field.
Then there are \(s_j\in\Gamma(U_j,\mathcal O_Z)\) with
\(n_{jj'}=s_j-s_{j'}\).
No cover or connectedness hypothesis on the index set
is needed.
\end{srlemma}
\begin{proof}
Every two nonempty opens of an integral scheme intersect,
and restriction of a section on a nonempty open to the
function field is injective.
For nonempty \(U_j\), form in \(K_Z\)
\[
                         r_j=r+n_{j j_0}.
\]
By (FC9) and characteristic \(p\),
\(r_j^p=h_{j_0}+h_j-h_{j_0}=h_j\).
At every point of \(U_j\), the element \(r_j\) is integral
over the local ring, since it solves the monic equation
\(T^p-h_j=0\).
Normality puts it in that local ring.
The local sections so obtained have the same function-field
germ, hence agree on overlaps and glue to a section \(s_j\).
This is the source's \srcid{exists\_section\_of\_pow} step.
For an empty \(U_j\), take its unique section.

On a nonempty \(U_j\cap U_{j'}\), the images in \(K_Z\) satisfy
\[
                    (s_j-s_{j'})^p=h_j-h_{j'}=n_{jj'}^p.
\]
Frobenius is injective in the field \(K_Z\), so their
function-field images are equal.
Injectivity of restriction gives the claimed equality
of sections.  On an empty intersection it is automatic.
\end{proof}

\subsection{The first-order consequence of a \(p\)-fold generator}
\label{subsec:faithful-qam-core}

The following is the source's \srcid{QamCore.core\_firstOrder},
with its thickenings and transition functions written out.

\begin{srlemma}\label{lem:faithful-qam-first-order}
Let \(X\) be a scheme, \(G\to X\) a scheme over \(X\),
\(k\) a field of characteristic \(p>0\), and
\(x:\operatorname{Spec}k\to X\).
Assume \(Z=G_x\) is integral and normal.
Put
\(Y_m=Z\times_k\operatorname{Spec}k[\delta]/(\delta^{m+1})\),
with projection \(b_m:Y_m\to G\) and zero section
\(\iota_m:Z\to Y_m\).
Let \(L\) be an invertible sheaf on \(G\).
Suppose \(a:Y_p\to G\) satisfies
\(a\iota_p=b_p\iota_p:Z\to G\), and suppose
\[
                    (a^*L\otimes b_p^*L^{-1})^{\otimes p}
\]
has a generator \(C\) whose restriction to \(Z\) is \(1\)
under the induced cancellation.

Take any trivializing cover \(V_j\) of \(L\), and put
\(U_j=Z\cap V_j\).
The coefficient of \(\delta^p\) in the local expression
of \(C\) on \(Y_p|_{U_j}\) is denoted \(h_j\).
Assume there is \(j_0\) with \(U_{j_0}\ne\varnothing\)
and a \(p\)-th root of \(h_{j_0}\) in \(K_Z\).
Then, after restriction to \(Y_1\),
\(a^*L\otimes b_1^*L^{-1}\) has a global generator.
\end{srlemma}
\begin{proof}
A nilpotent thickening has the same underlying topological
space as its closed subscheme.
Hence \(a\) and \(b_p\) have the same point map, and both
pull back \(V_j\) to the open \(Y_p|_{U_j}\).
Sections there have unique expansions
\[
                     \sum_{\nu=0}^{p}c_\nu\delta^\nu,
                \qquad c_\nu\in\Gamma(U_j,\mathcal O_Z).
 \tag{FC10}\label{eq:faithful-thick-expansion}
\]
For an affine open this is the tensor-product basis
\(1,\delta,\ldots,\delta^p\) of
\(k[\delta]/(\delta^{p+1})\).
On arbitrary opens, affine uniqueness makes the
coefficients agree under restriction and thus glue.
This is exactly the affine-basis and restriction argument
of \srcid{PowerSeriesExpansionInDeltaUnique.sections\_thick}.

Write \(g_{jj'}\) for the transition units of \(L\).
The line \(a^*L\otimes b_p^*L^{-1}\) has transition ratio
\[
 R_{jj'}=\frac{a^*g_{jj'}}{b_p^*g_{jj'}}
        =1+\sum_{\nu=1}^p n_{jj',\nu}\delta^\nu.
 \tag{FC11}\label{eq:faithful-thick-transition}
\]
The constant is \(1\) because \(a\) and \(b_p\) agree on \(Z\).
The local expressions of the given generator are
\[
                C_j=1+\sum_{\nu=1}^p c_{j,\nu}\delta^\nu,
                 \qquad C_j=R_{jj'}^p C_{j'} .
 \tag{FC12}\label{eq:faithful-generator-overlap}
\]
The coefficient \(c_{j,p}\) is \(h_j\).

Here is the source's coefficient comparison, including
the case \(p=2\).
In characteristic \(p\), with \(\delta^{p+1}=0\),
\[
 R_{jj'}^p=1+n_{jj',1}^p\delta^p:
\]
every exponent \(\nu p\) with \(\nu\ge2\) is at least \(p+1\).
Multiplication by a series with constant coefficient \(1\)
leaves this new \(\delta^p\)-coefficient unchanged.
Uniqueness in (FC10), applied to (FC12), gives
\[
                            h_j=h_{j'}+n_{jj',1}^p.
 \tag{FC13}\label{eq:faithful-thick-top-coefficient}
\]
Apply Lemma~\ref{lem:faithful-cech-p-root}.
There are \(s_j\in\Gamma(U_j,\mathcal O_Z)\) such that
\(n_{jj',1}=s_j-s_{j'}\).

On \(Y_1|_{U_j}\) the elements
\[
                         t_j=1+\delta s_j
\]
are units, with inverses \(1-\delta s_j\).
Since \(\delta^2=0\), (FC11) becomes
\[
  R_{jj'}=1+\delta(s_j-s_{j'})
                 =(1+\delta s_j)(1-\delta s_{j'}).
\]
Thus \(t_j=R_{jj'}t_{j'}\), which is the generator
compatibility relation on the diagonal trivializing charts.
To see that this is a generator of the actual line
on its full pullback cover, use the following explicit
extension, as in \srcid{exists\_gen\_of\_units}.
On a mixed chart \(a^{-1}V_i\cap b_1^{-1}V_j\), put
\[
                         a^*g_{ij}\,t_j.
\]
The two cocycle relations for \(g\), together with
\(t_j=R_{jj'}t_{j'}\), verify compatibility on every
mixed-chart intersection.
All these expressions are units, so they form a global
nowhere vanishing generator.
\end{proof}

\subsection{Exact dependency ledger for the next steps}
\label{subsec:faithful-frobenius-initial-ledger}

The two lemmas just proved are bounded source lemmas.
They do not yet prove (FC3), (FC6), or (FC8).
The remaining source route is:
\[
 \begin{gathered}
 \srcid{CbFrob.first\_power\_char\_p}
 \quad+\quad
 \srcid{AlphaKernelDerivativeVanishes.germKill}\\
 \Longrightarrow
 \srcid{CocpMain.cocycle\_p}
 \Longrightarrow
 \srcid{CocpMain.tanpf\_inj}\\
 \Longrightarrow
 \srcid{FamilyEx.fib\_inj},\ \srcid{FamilyEx.family\_ex}\\
 \Longrightarrow
 \srcid{FamilyGivesConjugateLieData.family\_conj}
 \Longrightarrow
 \srcid{FamilyHodgeLatticeData.family\_conj\_H\_eq}\\
 \Longrightarrow
 \srcid{SplitHermitianGeneratorsExist.deg\_quot\_eq\_zero}.
 \end{gathered}
\]
The first arrow also uses the actual \(p\)-fold cubical
generator \srcid{Pf.pfold}, the actual germ Taylor expansion,
and the normal function-field criterion
\srcid{pderiv\_zero\_of\_jet\_pth\_power\_diff}.
The second requires (FC4).
The third uses geometric-to-tangent power freeness,
base change of actual germs, common algebraically closed
points, and local affine germ existence.
The exact-sequence arrows require both formal coordinate
and gauge transformation identities, residue injectivity,
and integral local matrix splittings.
The last arrow retains the local free \(O\)-basis hypothesis
in (FC8) and uses the source's finite-field idempotent and
ideal-filtration argument.

For reproducibility, the frozen source paths and blob hashes
for this chunk are:
\begin{itemize}
\item \srcid{src/Development/AutomorphicSpectrum/Definitions/GlobalGeometry/ToFirst.lean}:
\srcid{e71c62079d11a8a5cf77513bf553601f0db05e76}.
\item \srcid{src/Development/AutomorphicSpectrum/Definitions/CommutativeAlgebra/GermComp1ZeroPointLeft.lean}:
\srcid{dac7d8db6b3addd6f6f8c112a17f50defc1ee274}.
\item \srcid{src/Development/AbelianVarietyTheory/Proofs/UnramifiedDepth/ExistsSectionOfPow.lean}:
\srcid{68fee6bc061158ef86e3065c1ad31d657447ac76}.
\item \srcid{src/Development/AutomorphicSpectrum/Proofs/GeometryProofs/ArtinSchreierCocycleRoot.lean}:
the wrapper invokes the preceding normal-stalk construction.
\item \srcid{src/Development/WeilBound/CoreStatements/HodgeTatePoints/QamCore.lean}:
\srcid{485616b11de8936fc073cf07462deb9f84fa455e}.
\item \srcid{src/Development/AutomorphicSpectrum/Proofs/GeometryProofs/PowerSeriesExpansionInDeltaUnique.lean}:
\srcid{3c57e495af93702da7f43bb550275016bfc15eec}.
\item \srcid{src/Development/AutomorphicSpectrum/Proofs/GeometryProofs/TruncatedExpansionPthPowerCoeff.lean}:
\srcid{fa9613feb577badff61581e9dae49879ec6e7149}.
\item \srcid{src/Development/AutomorphicSpectrum/Proofs/GeometryProofs/LineBundleCocyclePullbackTransform.lean}:
\srcid{a181938c65d1c8edc97bf19c8a49aabd1eef075e}.
\end{itemize}
The global-model construction, full polarized Mumford
uniformization, the free integral \(O\)-basis in (FC8), and
the analytic flat frame's canonical extension are separate
joins.  No conclusion about those joins is inferred here.

\subsection{The two formal Frobenius factorizations}
\label{subsec:faithful-frobenius-formal-factorizations}

This subsection translates the coefficient and Jacobian argument of
\srcid{CbFrob}, together with
\srcid{AlphaKernelDerivativeVanishes}.
It applies over every commutative ring \(R\) of characteristic
the prime \(p\), to the actual formal biextension datum
\(B:\mathrm{FB}(R,d,O)\).
No perfectness, reducedness or field assumption on \(R\)
is used in these formal lemmas.

Write \(F(X,Y)\) for its commutative formal group law,
\([m](X)\) for its specified multiplication by \(m\), and
\(\phi_1(X,Y,Z)\), \(\phi_2(X,Y,Z)\) for the two partial
multiplication series.
All substituted points below have zero constant term.
The group law has \(F(X,0)=F(0,X)=X\).
The partial multipliers are \(1\) whenever one of their
three arguments is zero.
The identities used here, which are fields of the
source's datum, include
\[
 \begin{split}
 \phi_1(X,Y,Z)&=\phi_1(Y,X,Z),\\
 \phi_1(X,Y,W)\phi_1(F(X,Y),Z,W)
   &=\phi_1(Y,Z,W)\phi_1(X,F(Y,Z),W),\\
 \phi_1(X,Y,Z)\phi_1(X,Y,W)\phi_2(F(X,Y),Z,W)
   &=\phi_2(X,Z,W)\phi_2(Y,Z,W)
                                  \phi_1(X,Y,F(Z,W)).
 \end{split}
 \tag{FP1}\label{eq:faithful-partial-laws}
\]
The second partial associativity identity is also retained
in the datum, though it is not needed for the first
Jacobian cancellation below.

Put
\[
                P_m(X,Z)=\prod_{\nu=0}^{m-1}
                                  \phi_1([\nu]X,X,Z).
 \tag{FP2}\label{eq:faithful-Pm}
\]
The empty product is \(1\).
The formal module laws give \([0]X=0\) and
\([m+1]X=F([m]X,X)\).

\begin{srlemma}\label{lem:faithful-Pm-identities}
For every \(m\ge0\),
\[
 \begin{split}
 P_m(F(X,Y),Z)\phi_1(X,Y,Z)^m
       &=P_m(X,Z)P_m(Y,Z)
                                \phi_1([m]X,[m]Y,Z),\\
 P_m(X,Y)P_m(X,Z)\phi_2([m]X,Y,Z)
       &=\phi_2(X,Y,Z)^m P_m(X,F(Y,Z)).
 \end{split}
 \tag{FP3}\label{eq:faithful-Pm-identities}
\]
\end{srlemma}
\begin{proof}
These are the source's two inductions on \(m\).
For the first, partial associativity and commutativity give
the four-term rearrangement
\[
 \begin{split}
 &\phi_1(A,B,Z)\phi_1(C,D,Z)
                              \phi_1(F(A,B),F(C,D),Z)\\
 &\hspace{8mm}=
 \phi_1(A,C,Z)\phi_1(B,D,Z)
                              \phi_1(F(A,C),F(B,D),Z).
 \end{split}
 \tag{FP4}\label{eq:faithful-four-term}
\]
To obtain it, compare the two orders of adjoining \(B,C,D\)
to \(A\) by partial associativity, and use commutativity to
interchange \(B,C\).
At \(m=0\), the first identity in (FP3) is the axes identity.
For the step \(m\mapsto m+1\), expand
\[
          P_{m+1}(X,Z)=P_m(X,Z)\phi_1([m]X,X,Z).
\]
Use (FP4) with
\((A,B,C,D)=([m]X,[m]Y,X,Y)\), the induction hypothesis,
and \([m]F(X,Y)=F([m]X,[m]Y)\).
The resulting products are precisely the first identity
at \(m+1\).

For the second identity, the case \(m=0\) is
\(\phi_2(0,Y,Z)=1\).
In the interchange identity (FP1), substitute
\(([m]X,X,Y,Z)\).
Multiply by \(P_m(X,Y)P_m(X,Z)\).
The left side becomes the expanded \(m+1\) left side.
On the right, the induction hypothesis replaces
\[
 P_m(X,Y)P_m(X,Z)\phi_2([m]X,Y,Z)
\]
by \(\phi_2(X,Y,Z)^mP_m(X,F(Y,Z))\).
The remaining factor is
\(\phi_2(X,Y,Z)\phi_1([m]X,X,F(Y,Z))\).
This gives the second identity at \(m+1\).
Every factor with constant coefficient \(1\) is a unit,
so the same computation can equivalently be read by
cancelling those factors, as in the source's polynomial
identity verification.
\end{proof}

\begin{srlemma}\label{lem:faithful-formal-Frobenius-factor}
There are unique formal series \(V(X)\) and \(g(X,Y)\)
such that
\[
 \begin{gathered}
 [p](X)=V(X_1^p,\ldots,X_d^p),\qquad V(0)=0,\\
 P_p(X,Y)=g(X_1^p,\ldots,X_d^p,Y),\\
                         g(0,Y)=g(X,0)=1.
 \end{gathered}
 \tag{FP5}\label{eq:faithful-Frobenius-factorizations}
\]
Moreover
\[
                       DV(0)=H,\qquad
                     \operatorname{bil}_{1,1}(g)=A,
 \tag{FP6}\label{eq:faithful-Frobenius-coefficients}
\]
for the matrices in \eqref{eq:faithful-alpha-beta}.
\end{srlemma}
\begin{proof}
We give the source's derivative argument.
Let
\[
                    J(X)_{ij}=
                           \partial_{Y_j}F_i(X,Y)|_{Y=0}.
\]
Its constant matrix is \(I_d\).
Thus \(\det J(X)\) has constant coefficient \(1\), and
\(J(X)\) is invertible over \(R[[X]]\).

Differentiate the formal homomorphism identity
\[
              [p](F(X,Y))=F([p]X,[p]Y)
\]
with respect to \(Y\) and set \(Y=0\).
Its right side has zero derivative because
\(D[p](0)=pI_d=0\).
Hence \(D[p](X)J(X)=0\); invertibility of \(J\) gives
every partial derivative of every component of \([p]\)
equal to zero.

Next, both derivatives in the first two arguments of
\(\phi_1\), evaluated at \((0,0,Z)\), vanish.
Indeed the corresponding axes restrictions are identically
\(1\).  Differentiating the finite product (FP2) at \(X=0\)
therefore gives
\[
                             D_XP_m(0,Z)=0
                                  \quad\text{for every }m.
 \tag{FP7}\label{eq:faithful-Pm-first-derivative}
\]
Differentiate the first identity in (FP3) for \(m=p\)
with respect to \(Y\), and set \(Y=0\).
The derivative of \(\phi_1(X,Y,Z)^p\) is zero in
characteristic \(p\).
On the other side, the derivative of \(P_p(Y,Z)\)
is zero by (FP7), and the derivative through \([p]Y\)
is zero by \(D[p](0)=0\).
The surviving identity is
\[
                            D_XP_p(X,Z)J(X)=0.
\]
Cancel \(J(X)\).
Thus every derivative of \(P_p\) in an \(X\)-variable
vanishes.

For a scalar series \(f=\sum_\nu c_\nu X^\nu\), this
derivative condition means that \(c_\nu=0\) unless every
\(\nu_i\) is divisible by \(p\).
If \(p\nmid\nu_i\), multiplication by \(\nu_i\) is
invertible in every ring of characteristic \(p\),
by Bezout's identity.
Consequently the only surviving exponents are \(p\)-multiples.
Reindexing coefficients constructs the unique series
after substitution \(X\mapsto X^{(p)}\).
This argument applies componentwise to \([p]\), and
only in the first block of variables to \(P_p\).
It proves (FP5).
The axes identities descend because that substitution
is coefficientwise injective.
The coefficient of \(X_j\) in \(V_i\) is the coefficient
of \(X_j^p\) in \([p]_i\), and the coefficient of \(X_iY_j\)
in \(g\) is that of \(X_i^pY_j\) in \(P_p\).
These are exactly (FP6).
\end{proof}

\subsection{The derivative killed by the kernel of \(\alpha\)}
\label{subsec:faithful-alpha-kernel-derivative}

For a series \(q\), let \(q^{[p]}\) mean applying
\(c\mapsto c^p\) to every coefficient.
This is different from substituting \(p\)-th powers
for its variables.

\begin{srlemma}\label{lem:faithful-alpha-kernel-derivative}
With the formal hypotheses above, let \(u\in R^d\) satisfy
\(\alpha_B(u)=0\), and let \(g\) be the series in (FP5).
Set
\[
                       D_u(Y)=
                       \sum_i u_i\,\partial_{X_i}g(0,Y).
 \tag{FP8}\label{eq:faithful-Du}
\]
Then every formal partial derivative of \(D_u(Y)\) is zero.
\end{srlemma}
\begin{proof}
The source first rewrites the second identity of (FP3)
at \(m=p\).
Substitute \([p]X=V(X^{(p)})\) and
\(P_p(X,Y)=g(X^{(p)},Y)\).
The coefficient Frobenius identity
\[
                 \phi_2(X,Y,Z)^p
                   =\phi_2^{[p]}(X^{(p)},Y^{(p)},Z^{(p)})
\]
holds term by term.
Substitution in the first variable block by \(p\)-th
powers is injective, so it can be cancelled.
One obtains the source's identity
\[
 \begin{split}
 g(X,Y)g(X,Z)\phi_2(V(X),Y,Z)
  =\phi_2^{[p]}(X,Y^{(p)},Z^{(p)})\,g(X,F(Y,Z)).
 \end{split}
 \tag{FP9}\label{eq:faithful-germKill-identity}
\]
This cancellation takes place in a power-series ring
and does not require Frobenius on \(R\) to be injective.

Apply \(\sum_i u_i\partial_{X_i}\), then set \(X=0\).
The axes identities make all undifferentiated factors \(1\).
The term from differentiating \(V(X)\) vanishes because
\(DV(0)u=Hu=0\), which is the first component of
\(\alpha_B(u)=0\).
Thus
\[
 D_u(Y)+D_u(Z)
   =E_u(Y^{(p)},Z^{(p)})+D_u(F(Y,Z)),
 \tag{FP10}\label{eq:faithful-germKill-additive}
\]
where
\[
             E_u(Y,Z)=
                 \sum_i u_i\partial_{X_i}
                                      \phi_2^{[p]}(0,Y,Z).
\]
Differentiate (FP10) with respect to \(Z_j\) and set \(Z=0\).
The derivative of \(E_u(Y^{(p)},Z^{(p)})\) is zero.
The chain rule therefore gives, for every \(j\),
\[
   \sum_k\partial_{Y_k}D_u(Y)\,J(Y)_{kj}
                                =\partial_{Y_j}D_u(0).
 \tag{FP11}\label{eq:faithful-germKill-Jacobian}
\]
By (FP6),
\[
                         \nabla D_u(0)=A^tu=0,
\]
the second component of \(\alpha_B(u)=0\).
The same Jacobian \(J(Y)\) used above is invertible.
Multiplying the row in (FP11) by its inverse proves
\(\partial_{Y_k}D_u=0\) for every \(k\).
This is the source's \srcid{germKill}; in particular no
assumption that \(\alpha\) is already injective has entered.
\end{proof}

\paragraph{Frozen-source scope and provenance.}
The lemmas above prove the formal-series part needed for
\srcid{CocpMain.cocycle\_p}.
They do not yet construct its global line generator or
deduce \(\mathrm{TanPF}\).
Their proof follows the same four steps as the frozen
source: the four-term induction, Jacobian cancellation,
coefficient reindexing, and the second Jacobian cancellation
for \(D_u\).
The source paths and checked blob hashes are
\[
\begin{array}{ll}
 \srcid{FirstPowerEqProdAt.lean}&
 \srcid{a8017b9ca8fc5ef04275c554343b57c74aabc657},\\
 \srcid{CbFrob.lean}&
 \srcid{575199c90bc9fec8a2509cc9f12c04ef6b492c3d},\\
 \srcid{CoeffEvPowOne.lean}&
 \srcid{01e43128ed3127558875f327af3d998e289c59b0},\\
 \srcid{AlphaKernelDerivativeVanishes.lean}&
 \srcid{5042e8658e3c7cf3417b9ebe62af86336fa43f0f}.
\end{array}
\]
The first three lie under
\srcid{src/Development/WeilBound/CoreStatements/HodgeTatePoints/};
the last lies under
\srcid{src/Development/AutomorphicSpectrum/Proofs/AlgebraProofs/}.
The formal data and identities are those of
\srcid{src/ReciprocityTheory/ReductiveGroups/ArithmeticGeometry/BirationalGeometry/ToroidalCompactification/CocycleRightZeroThd.lean}.
All are read at \srcid{alpoge-source-7cbf7a0c}.

\subsection{The source's Taylor map and function-field \(p\)-th root}
\label{subsec:faithful-taylor-function-field-root}

This is the route used in
\srcid{JetOfMono.taylor} and \srcid{CgTaylor.fibre\_taylor}.
In particular, the function-field root is obtained through
Kahler differentials and the source's derivation criterion,
not by an unstated comparison with a completed de Rham module.

\begin{srlemma}\label{lem:faithful-field-derivation-root}
Let \(K\) be any field of characteristic the prime \(p\).
If \(f\in K\) satisfies \(D(f)=0\) for every
\(\mathbb Z\)-derivation \(D:K\to K\), then \(f=h^p\)
for some \(h\in K\).
\end{srlemma}
\begin{proof}
We reproduce the source's maximal-subfield argument.
Suppose \(f\notin K^p\).
By Zorn's lemma choose a maximal subfield \(M\subset K\)
which contains \(K^p\) and does not contain \(f\).
The union of a nonempty chain of such subfields has the
same properties, which verifies the chain condition.

For every \(z\notin M\), its minimal polynomial over \(M\) is
\[
                              X^p-z^p.
 \tag{FT1}\label{eq:faithful-pure-minpoly}
\]
Indeed this polynomial annihilates \(z\), and it is
irreducible: a \(p\)-th root in \(M\) of \(z^p\) would equal
\(z\), by injectivity of Frobenius in \(K\).
The prime-degree irreducibility criterion for
\(X^p-c\) therefore applies.
Thus \([M(z):M]=p\).

Maximality implies \(f\in M(z)\) for every \(z\notin M\):
otherwise adjoining \(z\) would produce a larger permitted
subfield.  Both \(M(f)\) and \(M(z)\) have degree \(p\),
so \(M(f)=M(z)\).
It follows that \(K=M(f)\).

The surjection \(M[X]\to K,\ X\mapsto f\), has kernel
\((X^p-f^p)\).
Polynomial differentiation preserves this ideal because
\((X^p-f^p)'=0\) in characteristic \(p\).
It descends to an \(M\)-derivation \(K\to K\) carrying
\(f\) to \(1\).
Restricting scalars gives a \(\mathbb Z\)-derivation,
contradicting the hypothesis.
\end{proof}

\begin{srlemma}\label{lem:faithful-local-series-root}
Let \(k\) be a perfect field of characteristic the prime \(p\).
Let \(A\) be a local integral \(k\)-algebra essentially of
finite type, with maximal ideal \(\mathfrak m\).
Let
\[
                   j:A\longrightarrow k[[X_1,\ldots,X_d]]
\]
be a \(k\)-algebra map, and put
\(\mathfrak q=(X_1,\ldots,X_d)\).
Assume the two precise approximation properties
\[
 \begin{split}
 j(a)\in\mathfrak q^m
             &\iff a\in\mathfrak m^m
                                  &&(a\in A,\ m\ge0),\\
 \forall m\ge0,\ \forall s\in k[[X]],\quad
 \exists a\in A:\quad j(a)-s&\in\mathfrak q^m .
 \end{split}
 \tag{FT2}\label{eq:faithful-Taylor-properties}
\]
If every formal partial derivative of \(j(f)\) is zero,
then \(f\) has a \(p\)-th root in \(\operatorname{Frac}A\).
\end{srlemma}
\begin{proof}
Vanishing derivatives imply that every nonzero coefficient
of \(j(f)\) has exponent divisible by \(p\) in each variable.
Since \(k\) is perfect, take the unique \(p\)-th root of
each such coefficient.  Reindexing the exponents gives
\(g\in k[[X]]\) with
\[
                                 g^p=j(f).
\]
This is the source's
\srcid{exists\_pow\_eq\_of\_pderiv\_eq\_zero}.

For each \(m\ge0\), use (FT2) to choose \(a\in A\) with
\(j(a)-g\in\mathfrak q^{m+1}\).
Then
\[
        j(a^p-f)=(j(a)-g)^p\in\mathfrak q^{m+1},
            \qquad a^p-f\in\mathfrak m^{m+1}.
 \tag{FT3}\label{eq:faithful-Taylor-p-approximation}
\]
Let \(\Omega=\Omega_{A/k}\) and \(d:A\to\Omega\).
The Leibniz rule gives
\[
                          d(\mathfrak m^{m+1})
                                     \subset\mathfrak m^m\Omega.
\]
This inclusion follows by induction: differentiate a
product of one element of \(\mathfrak m^m\) and one of
\(\mathfrak m\), and use additivity.
Since \(d(a^p)=0\), (FT3) puts \(df\) in
\(\mathfrak m^m\Omega\) for every \(m\).

The ring \(A\) is Noetherian and \(\Omega\) is a finite
\(A\)-module because \(A\) is essentially of finite type
over \(k\).
Krull intersection for this finite module over the local
ring gives \(\bigcap_m\mathfrak m^m\Omega=0\).
Therefore \(df=0\).

Put \(K=\operatorname{Frac}A\).
Every \(\mathbb Z\)-derivation \(D:K\to K\) kills the
image of \(k\): each \(c\in k\) is a \(p\)-th power in \(k\),
and the Leibniz formula gives \(D(c)=0\).
Its restriction to \(A\) is a \(k\)-derivation.
The universal property of \(\Omega_{A/k}\) and \(df=0\)
then give \(D(f)=0\).
Lemma~\ref{lem:faithful-field-derivation-root} supplies
the required \(p\)-th root in \(K\).
\end{proof}

\begin{srlemma}\label{lem:faithful-universal-jets-taylor}
Let \(k\) be a perfect field of characteristic the prime \(p\).
Let \(\pi:Z\to\operatorname{Spec}k\) be an integral scheme
locally of finite type, and let \(e\in Z(k)\).
Suppose there are compatible maps
\[
 x_m:\operatorname{Spec}J_m\longrightarrow Z,\qquad
 J_m=k[X_1,\ldots,X_d]/(X_1,\ldots,X_d)^{m+1},
\]
reducing to \(e\), with the following universal property:
for every augmented \(k\)-algebra \(B\to k\) whose
augmentation ideal has \((m+1)\)-st power zero,
every \(k\)-morphism \(\operatorname{Spec}B\to Z\)
reducing to \(e\) factors uniquely through \(x_m\),
compatibly with its augmentation.
For every open \(U\) containing \(e\) and every
\(f\in\Gamma(U,\mathcal O_Z)\), there is
\(F\in k[[X_1,\ldots,X_d]]\) whose truncation to \(J_m\)
is \(x_m^*f\), for every \(m\).
It satisfies
\[
 \begin{split}
 F=0&\Longrightarrow f=0,\\
 \partial_{X_i}F=0\ \text{for all }i
   &\Longrightarrow
       \exists r\in K_Z:\ r^p=f\text{ in }K_Z .
 \end{split}
 \tag{FT4}\label{eq:faithful-Taylor-function-field}
\]
\end{srlemma}
\begin{proof}
Let \(A=\mathcal O_{Z,e}\).
The residue field is \(k\).
Each \(J_m\) is local, with nilpotent augmentation ideal,
so \(x_m\) factors through the canonical monomorphism
\(\operatorname{Spec}A\to Z\).
This gives an augmented map \(\tau_m:A\to J_m\).
We verify the exact statements
\[
          \tau_m\text{ is surjective},\qquad
                         \ker\tau_m=\mathfrak m^{m+1}.
 \tag{FT5}\label{eq:faithful-stalk-jet-quotient}
\]
The inclusion \(\mathfrak m^{m+1}\subset\ker\tau_m\)
follows from the augmentation.
Apply universal jets to
\(\operatorname{Spec}(A/\mathfrak m^{m+1})\to Z\);
the result is a map
\(\psi:J_m\to A/\mathfrak m^{m+1}\).
The monomorphism \(\operatorname{Spec}A\to Z\) implies
\(\psi\tau_m\) is the quotient map.
If \(\tau'_m:A/\mathfrak m^{m+1}\to J_m\) is induced
by \(\tau_m\), uniqueness in the universal property for
\(\operatorname{Spec}J_m\to Z\) gives
\(\tau'_m\psi=\mathrm{id}\).
Surjectivity of the quotient map also gives
\(\psi\tau'_m=\mathrm{id}\).
Thus \(\tau'_m\) is an isomorphism, proving (FT5).
This is the source's \srcid{jet\_of\_mono} argument.

The \(\tau_m\) commute with truncation: the resulting
maps into \(\operatorname{Spec}A\) have the same composite
to \(Z\), so are equal.
A compatible sequence of truncated polynomials determines
a unique formal power series, coefficient by coefficient.
Applying this to \((\tau_m(a))_m\) constructs a
\(k\)-algebra map \(j:A\to k[[X]]\).
Its algebra laws follow by checking every truncation.
Equation (FT5) gives the first assertion in (FT2).
Surjectivity of each \(\tau_m\) gives the approximation
assertion there, by lifting the truncation of any series.

Apply \(j\) to the germ of \(f\) at \(e\).
This gives \(F\), with the required truncations.
If \(F=0\), (FT2) puts the germ of \(f\) in every
\(\mathfrak m^m\).
The local ring is Noetherian, so Krull intersection
makes that germ zero.
On an integral scheme the map from sections on a
nonempty open to any stalk in that open is injective;
hence \(f=0\).
For the derivative assertion, apply
Lemma~\ref{lem:faithful-local-series-root}, and transport
the resulting root along
\(\operatorname{Frac}(\mathcal O_{Z,e})\simeq K_Z\).
\end{proof}

\paragraph{Application to the actual semiabelian fibre.}
Let \(t,x,g\) have the geometric hypotheses in
\eqref{eq:faithful-fibre-inj-target}, with \(k\)
algebraically closed.
The fibre \(Z=G_x\) is smooth over \(k\) and geometrically
connected, hence integral.
Its local rings are regular and therefore normal.
The fibre is locally of finite type, and \(k\) is perfect.
The actual germ maps \(g.x_m\), after passage to \(Z\),
satisfy precisely the universal jet property above.
In the source this is obtained from
\srcid{CdUniv.germ\_univ}: the compatible formal coordinates
are the universal parameters, their differential is the
specified tangent basis, and inverse coordinates are unique.
Pullback uniqueness identifies these maps with the
\(\tau_m\) used in (FT5).
Their compatibility and zero-section equations are the
truncation and zero clauses of the actual germ.
Thus \srcid{CgTaylor.fibre\_taylor} supplies (FT4) on
every open containing the identity.

This last paragraph uses the source's actual
\srcid{GermAt} and its universal-coordinate theorem.
It is not an assertion that arbitrary formal group data
come from an algebraic fibre.
The proofs of the coordinate theorem and the formal
group's construction are retained as dependencies of the
actual germ construction; the field-root step above is
fully proved once that actual germ is supplied.

\paragraph{A finite-jet derivative test.}
The source also uses the following coefficient fact.
If \(F,f\in k[[X_1,\ldots,X_d]]\), all derivatives of
\(f\) are zero, and for every \(m\) there is \(v_m\in J_m\)
with
\[
                           F\bmod(X)^{m+1}
                                  =v_m^p-f\bmod(X)^{m+1},
 \tag{FT6}\label{eq:faithful-jet-derivative-test}
\]
then all derivatives of \(F\) are zero.
To verify a coefficient of \(\partial_{X_i}F\) indexed
by \(\nu\), use \(m=|\nu|+1\), lift \(v_m\) to a polynomial
\(P\), and compare the coefficient indexed by
\(\nu+\mathbf e_i\) in \(F+f-P^p\).
It vanishes because that expression is zero modulo
\((X)^{m+1}\).
The derivatives of both \(f\) and \(P^p\) are zero,
so this is the desired derivative coefficient of \(F\).
This reproduces
\srcid{PderivZeroOfJetPthPowerDiff.pderiv\_eq\_zero\_of\_jet\_pow\_sub}.

\paragraph{Frozen-source ledger.}
The field derivation proof is
\srcid{src/Development/AbelianVarietyTheory/Proofs/UnramifiedDepth/ExistsPthRootOfDerivationZero.lean}
(blob \srcid{2c23c05250df5f7899c3be0fafd48c516142fe1d}).
The local differential/Krull proof is
\srcid{src/Development/AbelianVarietyTheory/Proofs/UnramifiedDepth/ExistsPowEqFractionRingOfPderivEqZero.lean}
(blob \srcid{261542edc4e53a309f6d0f23d8b3e73c32243079}).
The actual stalk-jet and Taylor construction is
\srcid{src/Development/WeilBound/CoreStatements/HodgeTatePoints/JetOfMono.lean}
(blob \srcid{f760cf4a69b739bfffb5979fd5638dc545657642}),
followed by \srcid{CgTaylor.lean}
(blob \srcid{fe0e799a60a74b234b78760406bf9b8e8abf0e25}).
The coefficient test is
\srcid{src/Development/AutomorphicSpectrum/Proofs/GeometryProofs/PderivZeroOfJetPthPowerDiff.lean}
(blob \srcid{cf8670ce38f1132e93262f544a8f6b54bcaa90a9}).
All hashes were checked against the frozen baseline tag.

The next join is the source's global cubical \(p\)-fold
generator and its exact finite-jet value.
That join, rather than the purely formal derivative
calculation alone, is what permits applying
Lemma~\ref{lem:faithful-qam-first-order}.

\subsection{The actual cubical \(p\)-fold generator}
\label{subsec:faithful-pfold-generator}

We now translate \srcid{CaNfold.nfold\_law} and
\srcid{Coc.Pf.pfold}, including the generator's finite-jet
value.  The family \(t\), its actual cubical line \(L\),
and its actual germ \(g\) are those in
\eqref{eq:faithful-fibre-inj-target}.
At this step \(k\) is any field of characteristic the prime \(p\).

For points of the group write the group law additively.
Set
\[
 \mathcal B_L(u,z)
       =L_{u+z}\otimes L_u^{-1}\otimes L_z^{-1}\otimes L_0.
\]
The cubical structure of the given line supplies the
normalized multiplication comparison for \(\mathcal B_L\).
Multiplying its cube section at
\(([\nu]u,u,z)\), for \(0\le\nu<n\), gives a generator
\[
 \kappa_n\quad\hbox{of}\quad
            \mathcal B_L([n]u,z)\otimes
                                     \mathcal B_L(u,z)^{-n}.
 \tag{FG1}\label{eq:faithful-kappa-line}
\]
Indeed the consecutive factors
\(\mathcal B_L([\nu+1]u,z)/\mathcal B_L([\nu]u,z)\)
telescope.  The remaining factors are
\(\mathcal B_L(u,z)^{-n}\).
The factor at \(\nu=0\) and the restriction at \(u=0\)
are the cubical axes trivializations.
This is the source's \srcid{exists\_kappa}, obtained by
taking the same finite product of its cubical generators.

Write \(e(u,z)\) for the actual axes-normalized formal
frame of \(\mathcal B_L(u,z)\) in the germ.
Its first partial multiplication identity is exactly
\[
 \phi_1(u,v,z)
  =\gamma(u,v,z)\,
             \frac{e(u,z)e(v,z)}{e(u+v,z)},
\]
where \(\gamma\) denotes the cubical section in those
line frames.
Multiplying at \(([\nu]u,u,z)\) yields
\[
              P_n(u,z)
                 =\kappa_n(u,z)\,
                               \frac{e(u,z)^n}{e([n]u,z)}.
 \tag{FG2}\label{eq:faithful-kappa-frame}
\]
This is an identity on every actual finite jet.
It proves the generator correspondence in
\srcid{CaNfold.nfold\_law}; it does not define a new
generator from generic-fibre data.

Put \(R_p=k[\delta]/(\delta^{p+1})\).
Let \(y\in G_x(R_p)\) reduce to zero and satisfy \([p]y=0\).
There is a generator \(\tau\) of
\[
                              y^*L\otimes0^*L^{-1}
 \tag{FG3}\label{eq:faithful-tau-line}
\]
whose reduction is the cancellation generator \(1\).
To construct it as in the source, choose one chart of \(L\)
containing the common image of \(y\) and \(0\); the spectrum
of \(R_p\) has just one point.
Use the ratio of its two chart frames.
On an arbitrary mixed chart \((i,j)\), its expression is
\[
                       y^*g_{i b}\,(0^*g_{j b})^{-1}.
\]
The transition cocycle for \(L\) proves these expressions
are a generator, and on reduction \(y=0\) their ratio is \(1\).
This is \srcid{LineBundleSectionDifferenceGenerator.jetPair\_gen}.

On \(Y_p=G_x\times_k\operatorname{Spec}R_p\), let
\(a_y(z,\delta)=z+y(\delta)\) and \(b(z,\delta)=z\).
Since \([p]y=0\), the first factor in (FG1) is canonically
trivial.
Pull \(\kappa_p\) back along \((y,z)\), and pull \(\tau\)
back from \(\operatorname{Spec}R_p\).
The product
\[
                             C=\kappa_p(y,z)^{-1}\tau^p
 \tag{FG4}\label{eq:faithful-C-generator}
\]
is a generator of
\[
                   (a_y^*L\otimes b^*L^{-1})^{\otimes p}.
\]
The line identity is explicit:
\(\kappa_p^{-1}\) generates \(\mathcal B_L(y,z)^p\);
multiplication by \((L_yL_0^{-1})^p\) leaves
\((L_{y+z}L_z^{-1})^p\).
Thus this is the actual translation-difference line used
in \eqref{eq:faithful-tanpf}.
At \(\delta=0\), the axes normalization of \(\kappa_p\)
and the normalization of \(\tau\) both give \(1\);
hence \(C|_{G_x}=1\).
These are exactly the construction and normalization in
\srcid{FrobeniusTorsionGeneratorRelation.exists\_gen\_pfold}
and \srcid{FrobeniusGeneratorCorrespondence.red\_one}.

\begin{srlemma}\label{lem:faithful-pfold-jet-value}
Suppose \(y\) is the actual germ point with formal
coordinates \(w\delta=(w_1\delta,\ldots,w_d\delta)\).
Let \(u=(w_i^p)_i\), and let \(g_0\) be the series of
\eqref{eq:faithful-Frobenius-factorizations}.
Set \(D_u(Z)=\sum_i u_i\partial_{X_i}g_0(0,Z)\).
For every \(m\ge0\), put
\[
 B_m=k[Z_1,\ldots,Z_d]/(Z)^{m+1},\qquad
 Q_m=B_m[\delta]/(\delta^{p+1}).
\]
Choose a chart of \(L\) containing the identity.
On the pullback of that chart along the actual germ
test \((w\delta,Z)\), there is \(v_m\in B_m\) for which
the generator in (FG4) has the exact expression
\[
                  1+\delta^p\bigl(v_m^p-D_u\bmod(Z)^{m+1}\bigr).
 \tag{FG5}\label{eq:faithful-pfold-jet-value}
\]
The expression concerns the actual generator \(C\), with
the same chart used in its \(\delta\)-coefficient expansion.
\end{srlemma}
\begin{proof}
Use the source's substitution map
\[
 q_m:
 k[X,Z]/(X,Z)^{p+m+2}\longrightarrow Q_m,\qquad
                         X_i\mapsto w_i\delta,\quad Z_i\mapsto Z_i.
\]
It is defined because a monomial of total degree at least
\(p+m+2\) must have either \(\delta\)-degree at least \(p+1\)
or \(Z\)-degree at least \(m+1\).
All the finite jets of the actual germs and frames commute
with this map, by their truncation and substitution clauses.

Let \(\nu_m\) be the image under \(q_m\) of the unit whose
underlying series is \(P_p\).
The factorization \(P_p(X,Z)=g_0(X^{(p)},Z)\) and the
axes identity \(g_0(0,Z)=1\) give
\[
                \nu_m
                =1+\delta^p\bigl(D_u\bmod(Z)^{m+1}\bigr).
 \tag{FG6}\label{eq:faithful-nu-jet}
\]
Indeed every term of degree at least two in \(X^{(p)}\)
contains \(\delta^{2p}=0\).
The surviving linear terms are
\(\delta^p\sum_iw_i^p\partial_{X_i}g_0(0,Z)\).
This is precisely
\srcid{GermJetQuotientFormula.jq\_psJet\_ptPow}.

Let \(e_m\) be the actual chart expression of \(e(y,Z)\),
and let \(\tau_m\) be the expression of the pullback of
\(\tau\).
Their reductions under \(\delta\mapsto0\) are both \(1\).
For \(e_m\), this is the actual frame's first axes identity;
for \(\tau_m\), it is (FG3)'s normalization.
The source verifies these equalities on the chosen chart
before taking their product.
By (FG2), \([p]y=0\), and the other axes identity,
the chart value of \(\kappa_p(y,Z)\) is
\(\nu_m e_m^{-p}\).
Thus the chart expression of (FG4) is
\[
                              \nu_m^{-1}(e_m\tau_m)^p.
 \tag{FG7}\label{eq:faithful-C-local-product}
\]

For any \(b\in B_m[\delta]/(\delta^{p+1})\) with
\(b|_{\delta=0}=1\), write
\(b=1+\delta(c+\delta h)\).
Characteristic \(p\) and \(\delta^{2p}=0\) imply
\[
                              b^p=1+\delta^pc^p.
\]
This is the source's
\srcid{JetAugmentationOnePthPower.pow\_eq\_one\_add};
it requires no reducedness of \(B_m\).
Apply it to \(b=e_m\tau_m\), obtaining some \(v_m\in B_m\).
Finally, the ideal generated by \(\delta^p\) has square zero,
so (FG6)--(FG7) give
\[
 \begin{split}
 (1+\delta^pD_u)^{-1}(1+\delta^pv_m^p)
                         &=1+\delta^p(v_m^p-D_u).
 \end{split}
\]
This is (FG5).
All expressions are taken in the same actual chart.
Their identities on germs of sections imply equality
of the chart sections, which is the final stalkwise
verification in \srcid{Pf.pfold}.
\end{proof}

\subsection{From the kernel vector to a global first-order generator}
\label{subsec:faithful-cocp-main}

\begin{srproposition}\label{prop:faithful-cocycle-p}
Let \(E,H_E,O,t,x,g\) have the precise geometric and
formal-germ meanings in
\eqref{eq:faithful-fibre-inj-target}.
Assume \(k\) is algebraically closed of characteristic
the prime \(p\).
For \(w\in k^d\), put \(u=(w_i^p)_i\).
If \(\alpha_{g.B}(u)=0\), then the first-order actual point
\(y_1\), with coordinates \(w\delta\), satisfies
\[
                         t_{y_1}^*L\otimes L^{-1}
                              \text{ has a global generator}.
 \tag{FG8}\label{eq:faithful-cocycle-p-conclusion}
\]
This is the source's \srcid{CocpMain.cocycle\_p}.
\end{srproposition}
\begin{proof}
Let \(y\in G_x(R_p)\) be the actual germ point with coordinates
\(w\delta\); its reduction is zero.
By \([p](X)=V(X^{(p)})\), its \(p\)-multiple has coordinates
\(V(u\delta^p)\).
Since \(V(0)=0\) and \(\delta^{2p}=0\), these are
\(\delta^p Hu\), which vanish by \(\alpha(u)=0\).
Thus \([p]y=0\).
This is the source's \srcid{pt\_of\_vec} calculation using
the actual jet maps, not an assertion only about tangent vectors.

Construct \(C\) by (FG4).
It is a normalized global generator of the \(p\)-th power
translation-difference line.
Choose a chart \(V_{j_0}\) of \(L\) containing the identity.
Every actual finite germ jet has underlying point the
identity, so all these jets lie in \(U_{j_0}=G_x\cap V_{j_0}\).
This open is nonempty.
Write \(h_{j_0}\) for the coefficient of \(\delta^p\)
in the \(\delta\)-expansion of \(C\) on this chart.

The actual Taylor construction of
Lemma~\ref{lem:faithful-universal-jets-taylor} gives a series
\(F\) for \(h_{j_0}\).
Pullback commutes with the finite free expansion in
\(1,\delta,\ldots,\delta^p\).
Therefore its coefficient of \(\delta^p\), after the
\(m\)-th actual germ test, is the truncation of \(F\).
Comparing with (FG5) gives, for every \(m\),
\[
                  F\bmod(Z)^{m+1}
                       =v_m^p-D_u\bmod(Z)^{m+1}.
 \tag{FG9}\label{eq:faithful-C-top-Taylor}
\]
This is the source's \srcid{DeltaExpansionRestriction.coef\_pull}:
the commuting pullback square identifies the two coefficient
maps, and uniqueness of the \(\delta\)-expansion identifies
their top coefficients.

Lemma~\ref{lem:faithful-alpha-kernel-derivative} gives
\(\partial D_u=0\).
The finite-jet test \eqref{eq:faithful-jet-derivative-test},
applied to (FG9), gives \(\partial F=0\).
The Taylor/function-field result
\eqref{eq:faithful-Taylor-function-field} supplies a
\(p\)-th root of \(h_{j_0}\) in the function field of \(G_x\).
The fibre \(G_x\) is integral and normal, as established
in the actual Taylor construction.
Apply Lemma~\ref{lem:faithful-qam-first-order}.
It gives a global generator of the translation-difference
line restricted to \(Y_1\), which is exactly (FG8).
\end{proof}

\begin{srcorollary}\label{cor:faithful-tanpf-injectivity}
With the same actual data, if
\(\mathrm{TanPF}(t,x,p)\) holds, then
\(\alpha_{g.B}\) is injective.
\end{srcorollary}
\begin{proof}
Let \(\alpha(u)=0\).
Since \(k\) is algebraically closed, choose \(w_i\in k\)
with \(w_i^p=u_i\).
The proof above constructs \(y\in G_x(R_p)\) reducing to
zero, with \([p]y=0\), and gives a global generator for
the translation-difference line of \(y_1\).
The exact definition \eqref{eq:faithful-tanpf} implies
\(y_1=0\).
The first-order formal coordinates are unique and have
the given tangent basis, so \(y_1=0\) forces every \(w_i=0\).
Therefore \(u=0\).
This is \srcid{IsPullbackWhisker.tanpf\_inj\_of}, and
with the preceding proposition it is
\srcid{CocpMain.tanpf\_inj}.
\end{proof}

\paragraph{Source and remaining scope.}
The construction follows the frozen source's ordered
choices: the global cubical \(\kappa_p\), the normalized
section \(\tau\), their product \(C\), the chart values
\(e_m,\tau_m\), the quotient unit \(\nu_m\), and the
nilpotent identity (FG7).
It uses no new alternate construction.

The main source files, under
\srcid{src/Development/AutomorphicSpectrum/Proofs/GeometryProofs/},
are \srcid{LineBundleGeneratorCorrespondence.lean},
\srcid{FrobeniusTorsionGeneratorRelation.lean},
\srcid{FrobeniusGeneratorCorrespondence.lean},
\srcid{LineBundleSectionDifferenceGenerator.lean},
\srcid{GermJetUnitOnChart.lean},
\srcid{JetSectionGeneratorValue.lean},
\srcid{GermFrobeniusTorsionSection.lean},
\srcid{GermJetQuotientFormula.lean},
\srcid{JetAugmentationOnePthPower.lean},
\srcid{DeltaExpansionRestriction.lean}, and
\srcid{UnitsMulRelationRearrangeCommRing.lean}.
Their last assembly file has frozen blob
\srcid{a780e8172520dc0915889c567ce4977f9ad3b70e}.
The cubical induction and final injectivity are in
\srcid{src/Development/WeilBound/CoreStatements/HodgeTatePoints/CaNfold.lean}
(blob \srcid{f16af35a98aa6f8756679ee6de758eaedc2e8480})
and \srcid{CocpMain.lean}
(blob \srcid{7aa25c6f73a189836a204c10956ead5c10783308}).
All were checked against the frozen baseline.

This proves the \(\mathrm{TanPF}\)-to-injectivity step.
The subsequent \(\mathrm{GeomPF}\)-to-\(\mathrm{TanPF}\)
descent, arbitrary-residue-field transport,
construction of the actual family over the arithmetic
curve, and integral splitting/gluing of (FC6)
are separate subsequent steps.
It supplies no independent model, quotient action,
or local free \(O\)-basis.
